\begin{textsample}{NEG Dim 4 – global\_south\_2024\_11 – Score 1.00 – global\_south\_2024\_11/117574609.txt}  \label{ex:f4_neg_173}
' Go to the nearest police station ', says Lamola on South Africans assisting Israel ' s army in Gaza war South African citizens are prohibited from participating in any foreign armed conflict.
Palestinians inspect the damage after an Israeli raid in the Nur Shams camp on 17 December 2023.
Picture : Zain Jaafar / AFP International Relations and Cooperation Minister Ronald Lamola says citizens can report to the police if they have evidence of any South Africans joining the Israel ' s army amid the war in Gaza.
Reports indicate that some South Africans have enlisted in the Israel Defense Forces ( IDF ) to engage in the ongoing conflict with Hamas in Gaza and other occupied Palestinian territories.
Since the onset of the Israel-Gaza conflict on 7 October 2023, over 43,000 Palestinians and 1,139 Israelis have lost their lives.
Lamola on South Africans in the Israel-Gaza war Speaking at a media briefing on Tuesday, Lamola said anyone with evidence of South Africans serving in the IDF should report it to the authorities. @ @ @ @ @ @ @ @ @ @ to the police and the NPA National Prosecuting Authority and the law must take its course for any South African participating in the Israel Defence Force,"he said."If we have evidence, that ' s good and we can refer it to the relevant law enforcement agencies, but the individuals who have identified and have got evidence of South Africans in the IDF, by law and right in this country, they can open a case against that individual at the nearest police station,"Lamola said.
Cases opened against citizens participating in Israel-Gaza war In April, the Palestine Solidarity Campaign ( PSC ) filed a criminal complaint against Benjamin Rattle for serving in the IDF.
Rattle had posted on Instagram, expressing pride in his role with the IDF.
The PSC submitted a list of about 70 names to the NPA.
More recently, the Directorate for Priority Crime Investigation, known as the Hawks, opened an investigation into 22-year-old Aaron Bayhack for his alleged involvement in the Israel-Gaza conflict @ @ @ @ @ @ @ @ @ @ spokesperson Siphiwe Dlamini stated that the military was aware of Bayhack ' s case and will allow legal proceedings to follow due process.
South Africans prohibited from joining other armies According to Dirco, South African law prohibits mercenary activities and regulates citizens or permanent residents joining foreign armed forces.
Under the Regulation of the Foreign Military Assistance Act, anyone wanting to provide military assistance abroad, including in Israel, must apply to the National Conventional Arms Control Committee ( NCACC ).
The NCACC will then recommend to the Minister of Defence whether to approve or deny the request.
Joining the IDF without NCACC approval is illegal and could lead to prosecution, as these actions may contribute to international law violations.
Additionally, under the South African Citizenship Act, naturalised citizens risk losing their citizenship if they fight for a foreign country in a conflict that South Africa does not support.

% matched lemmas: 
\end{textsample}
